\documentclass[aspectratio=169]{beamer}
\usepackage[english]{babel}
\usepackage[utf8]{inputenc}
\usepackage[T1]{fontenc}
\usepackage{lmodern}
\usepackage{listings}
\usepackage{xcolor}
\usepackage{graphicx}
\usepackage{textcomp}
\DeclareUnicodeCharacter{2014}{---}
\DeclareUnicodeCharacter{2013}{--}
\DeclareUnicodeCharacter{2212}{-}
\DeclareUnicodeCharacter{2192}{\ensuremath{\rightarrow}}
\DeclareUnicodeCharacter{00B1}{\ensuremath{\pm}}
\DeclareUnicodeCharacter{00D7}{\ensuremath{\times}}
\DeclareUnicodeCharacter{00B7}{\ensuremath{\cdot}}
\DeclareUnicodeCharacter{20AC}{EUR}
\DeclareUnicodeCharacter{2070}{\textsuperscript{0}}
\DeclareUnicodeCharacter{00B9}{\textsuperscript{1}}
\DeclareUnicodeCharacter{00B2}{\textsuperscript{2}}
\DeclareUnicodeCharacter{00B3}{\textsuperscript{3}}
\DeclareUnicodeCharacter{2074}{\textsuperscript{4}}
\DeclareUnicodeCharacter{2075}{\textsuperscript{5}}
\DeclareUnicodeCharacter{2076}{\textsuperscript{6}}
\DeclareUnicodeCharacter{2077}{\textsuperscript{7}}
\DeclareUnicodeCharacter{2078}{\textsuperscript{8}}
\DeclareUnicodeCharacter{2079}{\textsuperscript{9}}
\DeclareUnicodeCharacter{2248}{\ensuremath{\approx}}
\DeclareUnicodeCharacter{2264}{\ensuremath{\le}}

% Colors for code
\definecolor{codebg}{RGB}{245,245,245}
\definecolor{codecomment}{RGB}{0,128,0}
\definecolor{codekeyword}{RGB}{0,0,255}
\definecolor{codestring}{RGB}{163,21,21}
\definecolor{bandcolor}{RGB}{200,50,50}

\lstdefinestyle{pythonstyle}{
    language=Python,
    backgroundcolor=\color{codebg},
    basicstyle=\ttfamily\scriptsize,
    keywordstyle=\color{codekeyword},
    commentstyle=\color{codecomment},
    stringstyle=\color{codestring},
    showstringspaces=false,
    breaklines=true,
    frame=single,
    rulecolor=\color{black},
    escapeinside={(*@}{@*)},
    moredelim=**[l][\color{bandcolor}\bfseries]{\#\ ---\ NEW},
    moredelim=**[l][\color{bandcolor}\bfseries]{\#\ ---\ CHANGED},
    literate=
      {á}{{\'a}}1 {é}{{\'e}}1 {í}{{\'i}}1 {ó}{{\'o}}1 {ú}{{\'u}}1
      {Á}{{\'A}}1 {É}{{\'E}}1 {Í}{{\'I}}1 {Ó}{{\'O}}1 {Ú}{{\'U}}1
      {ñ}{{\~n}}1 {Ñ}{{\~N}}1 {ü}{{\"u}}1 {Ü}{{\"U}}1
      {¿}{{?`}}1 {¡}{{!`}}1
      {—}{{---}}1 {–}{{--}}1 {−}{{-}}1
      {±}{{\ensuremath{\pm}}}1 {×}{{\ensuremath{\times}}}1
      {→}{{\ensuremath{\rightarrow}}}1
      {°}{{\textdegree}}1
      {·}{{\ensuremath{\cdot}}}1
      {€}{{\texteuro}}1
      {≈}{{\ensuremath{\approx}}}1
      {≤}{{\ensuremath{\le}}}1
      {⁰}{{\textsuperscript{0}}}1 {¹}{{\textsuperscript{1}}}1
      {²}{{\textsuperscript{2}}}1 {³}{{\textsuperscript{3}}}1
      {⁴}{{\textsuperscript{4}}}1 {⁵}{{\textsuperscript{5}}}1
      {⁶}{{\textsuperscript{6}}}1 {⁷}{{\textsuperscript{7}}}1
      {⁸}{{\textsuperscript{8}}}1 {⁹}{{\textsuperscript{9}}}1
}

\lstset{style=pythonstyle}

% Title slide macro: \lessontitle{Lesson Number}{Title}
\newcommand{\lessontitle}[2]{
    \title{Lesson #1: #2}
    \author{}
    \institute{Tec de Monterrey}
    \begin{frame}[plain]
        \titlepage
    \end{frame}
}

% Figure frame macro: \figureframe{Title}{Image path}
\newcommand{\figureframe}[2]{
    \begin{frame}{#1}
        \begin{center}
            \includegraphics[height=0.8\textheight,width=\textwidth,keepaspectratio]{#2}
        \end{center}
    \end{frame}
}


% Listings pulled from src/ with \lstinputlisting, so the slide is the file.
\lstset{
    basicstyle=\ttfamily\fontsize{8pt}{9.6pt}\selectfont,
    breaklines=true,
    breakatwhitespace=true,
    columns=fullflexible,
    keepspaces=true,
    xleftmargin=0.3em,
    framesep=2pt,
    aboveskip=0pt,
    belowskip=0pt
}

\newcommand{\resultframe}[3]{%
    \begin{frame}{#1}
        \begin{center}
            \includegraphics[height=0.62\textheight,width=\textwidth,keepaspectratio]{#2}
        \end{center}
        \vspace{0.25em}
        {\small #3\par}
    \end{frame}%
}

\date{}

\begin{document}

\lessontitle{01}{Your First Genetic Algorithm}

\section{This lesson}

\begin{frame}{This lesson}
The course introduction showed how a lesson is read and how to run the code. Today the first algorithm is built, one operator at a time.

\vspace{0.6em}
\begin{itemize}
    \item A population, not a point; selection discards, crossover recombines, mutation invents.
    \item The loop of those three, for a handful of generations.
    \item The same code, a different seed, a local hill: a GA is not a guarantee.
\end{itemize}

\vspace{0.4em}
{\small The main route develops the algorithm. A code-reading appendix explains
the Python, plotting, paths, random state, and validation idioms used around it.}
\end{frame}

% ================= first_example_01_the_landscape =================
\begin{frame}{Mathematical model}
\[
\max_{x\in[-10,10]} f(x),\qquad f(x)=\sin x-0.2|x|
\]
An individual stores a candidate \(x\); its fitness is \(f(x)\). A population
is a finite sample, so its champion is the \emph{best observed} value. The dense
grid used below is a numerical reference, not a proof of the exact maximizer.
\end{frame}

\section{The problem}

\begin{frame}{The problem}
Before writing any algorithm, look at what we are searching.

\[
    f(x) = \sin(x) - 0.2\,|x|, \qquad x \in [-10,\,10]
\]

Maximise. No derivatives. Nothing about the function except how to evaluate it.
\end{frame}

\begin{frame}[fragile]{objective()}
\lstinputlisting[firstline=33,lastline=33]{../src/first_example_01_the_landscape.py}
\lstinputlisting[firstline=49,lastline=51]{../src/first_example_01_the_landscape.py}
\end{frame}

\resultframe{The search space}{../figures/first_example_01_the_landscape.png}{%
    Largest value on the 400-point grid: \(x = +1.378\), \(f = +0.706\).
    Multiple peaks: a method that only climbs uphill stops at whichever one it starts on.}

% ================= first_example_02_random_population =================
\section{A population}

\begin{frame}{A population}
A genetic algorithm does not track one candidate solution, it tracks many.
Here we simply scatter ten of them at random and look at where they land.
\end{frame}

\begin{frame}[fragile]{(1a) Individual: state and fitness}
\lstset{basicstyle=\ttfamily\fontsize{6.5pt}{7.8pt}\selectfont}
\lstinputlisting[firstline=58,lastline=59]{../src/first_example_02_random_population.py}
\lstinputlisting[firstline=71,lastline=82]{../src/first_example_02_random_population.py}
\end{frame}

\begin{frame}[fragile]{(1b) Individual: access and representation}
\lstset{basicstyle=\ttfamily\fontsize{6.5pt}{7.8pt}\selectfont}
\lstinputlisting[firstline=83,lastline=105]{../src/first_example_02_random_population.py}
\end{frame}

\begin{frame}[fragile]{(2) create\_random()}
\lstinputlisting[firstline=109,lastline=110]{../src/first_example_02_random_population.py}
\lstinputlisting[firstline=123,lastline=125]{../src/first_example_02_random_population.py}
\end{frame}

\begin{frame}[fragile]{(3) the population}
\lstinputlisting[firstline=129,lastline=141]{../src/first_example_02_random_population.py}
\end{frame}

\resultframe{Ten blind guesses}{../figures/first_example_02_random_population.png}{%
    Best of the initial population: \(x = -0.323\), \(f = -0.382\).
    Nowhere near the optimum.}

% ================= first_example_03_selection =================
\section{Selection}

\begin{frame}{Selection}
The population from step 2 is random and mostly bad. Selection is the first
force of evolution: it decides who gets to reproduce. Nothing new is created
here --- we only choose, with repetition, from what already exists.
\end{frame}

\begin{frame}[fragile]{(1) select\_tournament()}
\lstinputlisting[firstline=121,lastline=122]{../src/first_example_03_selection.py}
\lstinputlisting[firstline=142,lastline=151]{../src/first_example_03_selection.py}
\end{frame}

\resultframe{Selection only copies}{../figures/first_example_03_selection.png}{%
    No new \(x\) value appears. 4 of 10 individuals were selected zero times.
    Those exact objects no longer have a selected copy from which to reproduce.}

% ================= first_example_04_crossover =================
\section{Crossover}

\begin{frame}{Crossover}
Selection alone can only copy. Crossover is the first operator that creates
values that were not in the population before, by combining two parents.
\end{frame}

\begin{frame}[fragile]{(1) clamp()}
\lstinputlisting[firstline=60,lastline=61]{../src/first_example_04_crossover.py}
\lstinputlisting[firstline=80,lastline=82]{../src/first_example_04_crossover.py}
\end{frame}

\begin{frame}[fragile]{(2) crossover\_blend()}
\lstinputlisting[firstline=129,lastline=130]{../src/first_example_04_crossover.py}
\lstinputlisting[firstline=155,lastline=163]{../src/first_example_04_crossover.py}
\end{frame}

\begin{frame}[fragile]{(3) crossover()}
\lstinputlisting[firstline=167,lastline=168]{../src/first_example_04_crossover.py}
\lstinputlisting[firstline=181,lastline=184]{../src/first_example_04_crossover.py}
\end{frame}

\resultframe{What \(\alpha > 0\) buys}{../figures/first_example_04_crossover.png}{%
    14 of 20 children fell outside the parents' interval \([-2, +3]\).
    With \(\alpha = 0\), crossover would interpolate only: no child could leave that interval.}

% ================= first_example_05_mutation =================
\section{Mutation}

\begin{frame}{Mutation}
Crossover has little reach after parent genes cluster closely on one hill.
Mutation adds an independent displacement, so it can still propose a distant
gene even when both parents are identical.
\end{frame}

\begin{frame}[fragile]{(1) mutate\_gaussian()}
\lstinputlisting[firstline=123,lastline=124]{../src/first_example_05_mutation.py}
\lstinputlisting[firstline=143,lastline=145]{../src/first_example_05_mutation.py}
\end{frame}

\begin{frame}[fragile]{(2) mutate()}
\lstinputlisting[firstline=149,lastline=150]{../src/first_example_05_mutation.py}
\lstinputlisting[firstline=164,lastline=167]{../src/first_example_05_mutation.py}
\end{frame}

\resultframe{Reach of a single mutation}{../figures/first_example_05_mutation.png}{%
    From a population trapped at \(x \approx -4.6\), 6 units from the global hill:\\
    \(\sigma = 1.0\): 0 of 2000 sampled mutations land on the global hill.\\
    \(\sigma = 3.0\): 110 of 2000 mutations (5.5\%) land on the global hill.}

% ================= first_example_06_the_full_loop =================
\section{The complete genetic algorithm}

\begin{frame}{The complete genetic algorithm}
\[
    \text{initialise} \;\to\;
    [\;\text{SELECT} \to \text{CROSSOVER} \to \text{MUTATE} \to \text{replace}\;] \times N
\]
Nothing below is new except the loop itself.
\end{frame}

\begin{frame}[fragile]{(1a) initialise and select}
\lstset{basicstyle=\ttfamily\fontsize{6.5pt}{7.8pt}\selectfont}
\lstinputlisting[firstline=200,lastline=218]{../src/first_example_06_the_full_loop.py}
\end{frame}

\begin{frame}[fragile]{(1b) vary, replace and record}
\lstset{basicstyle=\ttfamily\fontsize{6.5pt}{7.8pt}\selectfont}
\lstinputlisting[firstline=219,lastline=240]{../src/first_example_06_the_full_loop.py}
\end{frame}

\begin{frame}[fragile]{(2) the convergence plot}
\lstinputlisting[firstline=250,lastline=264]{../src/first_example_06_the_full_loop.py}
\end{frame}

\resultframe{Ten generations}{../figures/first_example_06_the_full_loop.png}{%
    Approximate solution: \(x = +1.372\), \(f = +0.706\), near the highest peak
    (400-point grid reference near \(x = +1.378\), \(f = +0.706\)).}

% ================= first_example_07_local_optimum =================
\section{When it fails}

\begin{frame}{When it fails}
\[
    \texttt{SEED = 52} \quad\longrightarrow\quad \texttt{SEED = 16}
\]
Same algorithm, same parameters, same number of generations. The seed changes
the complete random sequence: initialization, tournaments, crossover and
mutation. This run settles on a local maximum and never leaves it.

This is not a bug to be fixed. It is the honest behaviour of the method, and
it is the reason Lessons 03, 04, 05 and 07 exist.
\end{frame}

\begin{frame}[fragile]{SEED}
\lstinputlisting[firstline=36,lastline=38]{../src/first_example_07_local_optimum.py}
\end{frame}

\resultframe{Stuck on the local hill}{../figures/first_example_07_local_optimum.png}{%
    Solution: \(x = -4.417\), \(f = +0.073\).
    Grid reference \(x = +1.378\), \(f = +0.706\) --- never found.}

\section{Next}

\begin{frame}{Next}
Seed 52 reaches the highest hill; seed 16 settles at \(x = -4.417\), \(f = +0.073\) for this ten-generation run. That failure is the method.

\vspace{0.8em}
\textbf{Lesson 02 --- The Genetic Algorithm Flow}

The loop we just wrote is still a run of unnamed comments. Lesson 02 names the five phases. Same seed, same digits as step 6.
\end{frame}

\appendix
\section{Code-reading appendix}

\begin{frame}{How to use this appendix}
The lesson assumes basic variables, functions and loops. These frames decode the
Python and plotting constructs that support the genetic algorithm.

\begin{itemize}
    \item Read each source excerpt from top to bottom; comments state its purpose.
    \item Follow the data: inputs enter on the right of an assignment, results receive the name on the left.
    \item Plotting and file-output statements present evidence; they do not change the GA state.
    \item The guided notebook gives the same explanations beside every executable cell.
\end{itemize}
\end{frame}

\begin{frame}[fragile]{Step 1: imports, settings and a portable path}
\lstset{basicstyle=\ttfamily\fontsize{6.5pt}{7.8pt}\selectfont}
\lstinputlisting[firstline=14,lastline=30]{../src/first_example_01_the_landscape.py}
\vspace{0.4em}
{\small Imports load reusable definitions. Uppercase names are run settings.
\texttt{Path / name} joins paths portably; \texttt{\_\_file\_\_} locates this script.}
\end{frame}

\begin{frame}[fragile]{Step 1: arrays, indexing and the plot lifecycle}
\lstset{basicstyle=\ttfamily\fontsize{5.2pt}{6.2pt}\selectfont}
\lstinputlisting[firstline=54,lastline=82]{../src/first_example_01_the_landscape.py}
\vspace{0.2em}
{\small \texttt{linspace} makes aligned inputs; \texttt{argmax} returns an index.
Create the , add layers and labels, save it, then close it.}
\end{frame}

\begin{frame}[fragile]{Step 2: object state, a property and formatted output}
\lstset{basicstyle=\ttfamily\fontsize{4.7pt}{5.4pt}\selectfont}
\lstinputlisting[firstline=71,lastline=105]{../src/first_example_02_random_population.py}
\vspace{0.2em}
{\small \texttt{self} is the new object. Fitness is cached when it is built.
\texttt{@property} hides the one-item list; \texttt{:+.3f} prints a sign and three decimals.}
\end{frame}

\begin{frame}[fragile]{Step 2: comprehensions, loops and a key function}
\lstset{basicstyle=\ttfamily\fontsize{6pt}{7.2pt}\selectfont}
\lstinputlisting[firstline=129,lastline=142]{../src/first_example_02_random_population.py}
\vspace{0.4em}
{\small A comprehension builds the population; the loop prints one object at a time.
\texttt{key=lambda ...} tells \texttt{max} which attribute to compare.}
\end{frame}

\begin{frame}[fragile]{Step 3: references, identity and counts}
\lstset{basicstyle=\ttfamily\fontsize{4.4pt}{5.1pt}\selectfont}
\lstinputlisting[firstline=142,lastline=178]{../src/first_example_03_selection.py}
\vspace{0.2em}
{\small Selection appends references to existing objects. \texttt{id} distinguishes
objects; \texttt{Counter.get(key, 0)} also represents candidates with no copies.}
\end{frame}

\begin{frame}[fragile]{Step 4: tuples, extend, filtering and Boolean counts}
\lstset{basicstyle=\ttfamily\fontsize{4.8pt}{5.6pt}\selectfont}
\lstinputlisting[firstline=189,lastline=221]{../src/first_example_04_crossover.py}
\vspace{0.2em}
{\small A comma returns a pair and unpacking names it. \texttt{extend} keeps one flat
list. Chained comparisons classify children; a Boolean contributes 1 or 0 to a count.}
\end{frame}

\begin{frame}{Two meanings of alpha}
\begin{description}
    \item[\texttt{BLEND\_ALPHA}] A GA parameter. It expands the interval from which blend crossover can produce children.
    \item[\texttt{alpha=0.2}] A Matplotlib argument. It makes a plotted layer 20\% opaque.
\end{description}

The shared name comes from two different APIs. Changing one has no effect on the other.
\end{frame}

\begin{frame}[fragile]{Step 5: a finite experiment and its criterion}
\lstset{basicstyle=\ttfamily\fontsize{6pt}{7.2pt}\selectfont}
\lstinputlisting[firstline=203,lastline=218]{../src/first_example_05_mutation.py}
\vspace{0.4em}
{\small Each sigma starts from corresponding random draws. The test counts a
neighbourhood of a hill. Zero observed arrivals do not establish zero probability.}
\end{frame}

\begin{frame}[fragile]{Step 5: the Gaussian density}
\lstset{basicstyle=\ttfamily\fontsize{5pt}{5.9pt}\selectfont}
\lstinputlisting[firstline=225,lastline=254]{../src/first_example_05_mutation.py}
\vspace{0.2em}
{\small The function evaluates the normal-density formula element by element.
The three aligned arrays support the figure; they do not alter mutation.}
\end{frame}

\begin{frame}[fragile]{Step 5: two y-axes and a merged legend}
\lstset{basicstyle=\ttfamily\fontsize{4.6pt}{5.4pt}\selectfont}
\lstinputlisting[firstline=256,lastline=291]{../src/first_example_05_mutation.py}
\vspace{0.2em}
{\small \texttt{twinx} shares x but gives density a different y scale. Each axis
owns its legend entries, so the two handle and label lists are joined explicitly.}
\end{frame}

\begin{frame}[fragile]{Step 6: slicing, zip and a flat offspring list}
\lstset{basicstyle=\ttfamily\fontsize{5.8pt}{7pt}\selectfont}
\lstinputlisting[firstline=213,lastline=226]{../src/first_example_06_the_full_loop.py}
\vspace{0.4em}
{\small \texttt{[::2]} and \texttt{[1::2]} take alternating slots. \texttt{zip}
aligns adjacent parents. \texttt{extend} adds both returned children separately.}
\end{frame}

\begin{frame}[fragile]{Step 6: conditional comprehension and replacement}
\lstset{basicstyle=\ttfamily\fontsize{6pt}{7.2pt}\selectfont}
\lstinputlisting[firstline=227,lastline=240]{../src/first_example_06_the_full_loop.py}
\vspace{0.4em}
{\small One random decision is made per candidate. Assignment replaces the list;
\texttt{append} then adds one scalar to history. No elitism protects the old best.}
\end{frame}

\begin{frame}[fragile]{Step 7: what the seed changes}
\lstset{basicstyle=\ttfamily\fontsize{6.5pt}{7.8pt}\selectfont}
\lstinputlisting[firstline=207,lastline=217]{../src/first_example_07_local_optimum.py}
\vspace{0.4em}
{\small \texttt{seed} fixes every subsequent pseudorandom draw. It changes more
than the initial population, so the two seeds compare complete random sequences.}
\end{frame}

\begin{frame}{Code-reading checklist}
For any unfamiliar chunk, ask:
\begin{enumerate}
    \item Which names already exist, and what type of value does each hold?
    \item Which expressions produce new values, objects, lists or figures?
    \item Does the chunk modify GA state, or only report and visualize it?
    \item Which random draw occurs here, and has the seed been reset?
    \item What invariant should be true afterward: size, bounds, cached fitness or reproducibility?
\end{enumerate}
\end{frame}
\end{document}
